\chapter{Bài toán, mục tiêu và phạm vi}
\label{ch:problem}

\section{Bối cảnh ứng dụng và động cơ nghiên cứu}
\label{sec:context}

Một phân xưởng có thể đồng thời chứa nhiều người, nhiều máy và nhiều khu vực với mức rủi ro khác nhau. Cùng một tư thế có thể mang ý nghĩa khác nhau tùy theo vị trí và vật thể xung quanh: cúi đầu có thể là kiểm tra sản phẩm, nhìn điện thoại hoặc bắt đầu ngủ gục; đưa tay ra phía trước có thể là thao tác đúng quy trình hoặc chạm vào vùng nguy hiểm. Vì vậy, bài toán phải được đặt ở mức \emph{từng người} và \emph{từng thời điểm}, không phải phân loại ảnh tĩnh hay gán một nhãn cho cả đoạn video.

Luận văn tập trung vào bốn hành vi mục tiêu: sử dụng điện thoại (\path{use_mobilephone}), ngủ trên bàn (\path{sleep_on_table}), gác chân lên bàn (\path{feet_on_table}) và hành vi nguy hiểm (\path{dangerous_action}). Thực nghiệm được tiến hành trên bộ dữ liệu riêng thu từ camera của phân xưởng, gán nhãn theo từng người tại các khung hình khóa (\cref{sec:factory-data}), và trên bộ dữ liệu công khai AVA~v2.2 \cite{gu2018ava}, vốn có cùng cấu trúc bài toán và là chuẩn đánh giá phổ biến nhất của STAD.

\section{Phát biểu bài toán}
\label{sec:formulation}

Theo định dạng AVA \cite{gu2018ava}, chú thích được gán tại các \emph{khung hình khóa} (keyframe). Tại mỗi khung hình khóa, mỗi người được mô tả bằng một hộp bao $b = (x_1, y_1, x_2, y_2)$ và một tập nhãn hành vi $\mathcal{Y} \subseteq \{1, \dots, K\}$; một người có thể mang nhiều nhãn đồng thời. Cho đoạn video $V \in \R^{3 \times T_0 \times H_0 \times W_0}$ lấy đối xứng quanh khung hình khóa và $N$ hộp người đề xuất $\{b_k\}_{k=1}^{N}$, mô hình cần trả ra với mỗi hộp $b_k$ một vectơ xác suất đa nhãn $\mat{p}_k \in [0,1]^K$, trong đó $p_{k,c}$ là xác suất người trong hộp $b_k$ đang thực hiện hành vi $c$. Với bộ dữ liệu phân xưởng $K = 4$; với AVA $K = 80$.

Pipeline hai giai đoạn phổ biến \cite{gu2018ava,feichtenhofer2019slowfast,mmaction2} thực hiện chuỗi biến đổi
\begin{equation}
\begin{aligned}
V &\xrightarrow{\text{backbone}} F \in \R^{C \times T \times H \times W}
  \xrightarrow{\text{RoIAlign}(b_k)} r_k \in \R^{C \times 1 \times h \times w}\\
  &\xrightarrow{\text{head}} \mat{p}_k = \sigma(\mat{W} \Pool(r_k) + \mat{b}),
\end{aligned}
\label{eq:pipeline}
\end{equation}
trong đó $\sigma$ là hàm sigmoid áp dụng độc lập trên từng lớp. Luận văn chèn một toán tử thời gian $\Phi$ vào giữa backbone và RoIAlign, $F \mapsto \Phi(F)$, để mỗi vị trí không gian của $F$ được làm giàu bằng thông tin dọc trục thời gian trước khi vùng người được cắt ra.

\section{Độ đo đánh giá}
\label{sec:metric}

Chất lượng phát hiện được đánh giá bằng \emph{frame-mAP} tại ngưỡng IoU~0,5 theo giao thức chính thức của AVA \cite{gu2018ava}. Với mỗi lớp $c$, các dự đoán được sắp xếp theo điểm $p_{k,c}$ giảm dần; một dự đoán là đúng nếu hộp của nó có IoU $\geq 0{,}5$ với một hộp nhãn thật chưa được ghép mang nhãn $c$. Khi đó
\begin{equation}
\mathrm{AP}_c = \int_0^1 P_c(R)\,\mathrm{d}R, \qquad
\mAP = \frac{1}{|\mathcal{C}_{\text{eval}}|} \sum_{c \in \mathcal{C}_{\text{eval}}} \mathrm{AP}_c,
\label{eq:map}
\end{equation}
với $\mathcal{C}_{\text{eval}}$ là bốn lớp mục tiêu trên dữ liệu phân xưởng và 60 lớp được đánh giá của AVA~v2.2. Bộ đánh giá \code{AVAMetric} của MMAction2 \cite{mmaction2} cài đặt giao thức này và báo cáo thêm mAP theo ba nhóm hành vi của AVA. Precision và Recall tại một ngưỡng cố định phụ thuộc mạnh vào hiệu chỉnh ngưỡng nên chỉ là chỉ số phụ; tuy vậy, chúng phản ánh trực tiếp hành vi của hệ cảnh báo tại một điểm vận hành, nên luận văn báo cáo hai chỉ số này tại ngưỡng 0,5, cùng tốc độ xử lý và bộ nhớ GPU.

\section{Đặc thù của các hành vi mục tiêu}
\label{sec:behaviours}

Bốn hành vi mục tiêu khác nhau rõ rệt về thang thời gian và tín hiệu nhận dạng (\cref{tab:behaviours}).

\begin{table}[htbp]
\centering
\caption{Đặc thù của bốn hành vi mục tiêu và yêu cầu kéo theo}
\label{tab:behaviours}
\begin{tabularx}{\textwidth}{lXX}
\toprule
\textbf{Hành vi} & \textbf{Đặc điểm tín hiệu} & \textbf{Yêu cầu đối với mô hình}\\
\midrule
Sử dụng điện thoại & Vật thể nhỏ, thường bị bàn tay che khuất; pha đưa điện thoại lên diễn ra nhanh & Giữ chi tiết không gian ở độ phân giải cao; nhận ra biến đổi ngắn hạn giữa các khung liền kề\\
Ngủ trên bàn & Chuyển động rất nhỏ, kéo dài nhiều giây; quan hệ đầu--thân với mặt bàn & Tích lũy trạng thái qua toàn bộ đoạn video; ngữ cảnh cảnh vật quanh người\\
Gác chân lên bàn & Gần như tĩnh; cần quan sát toàn thân, bàn chân và bề mặt đỡ & Trường tiếp nhận không gian rộng; xác nhận trạng thái bền vững theo thời gian\\
Hành vi nguy hiểm & Thay đổi nhanh, đa dạng, phụ thuộc thiết bị xung quanh & Giữ chi tiết cục bộ theo thời gian; quan hệ với thiết bị trong cảnh\\
\bottomrule
\end{tabularx}
\end{table}

Bốn hành vi tạo ra hai thang thời gian đối lập: \path{sleep_on_table} và \path{feet_on_table} cần quan sát dài để phân biệt trạng thái bền vững với tư thế tạm thời, trong khi \path{dangerous_action} và pha chuyển tiếp của \path{use_mobilephone} có thể chỉ kéo dài vài khung hình. Một toán tử thời gian với trường tiếp nhận cố định buộc phải thỏa hiệp giữa hai yêu cầu này. Đây là động cơ trực tiếp cho thiết kế hai nhánh của Hybrid FGTN-Mamba (\cref{ch:method}): một nhánh truyền trạng thái trên toàn bộ trục thời gian, một nhánh so khớp trực tiếp các khung hình liền kề, hợp nhất bằng cổng học được thay vì tỷ lệ cố định.

Nhận xét trên cũng áp dụng cho AVA: nhóm \emph{chuyển động của người} (đi, đứng, ngồi, nằm) gồm các trạng thái bền vững; nhóm \emph{thao tác với vật thể} (cầm, uống, viết, nghe điện thoại) phụ thuộc vật thể nhỏ và các pha ngắn; nhóm \emph{tương tác giữa người} (nói chuyện, nghe, nhìn người khác) phụ thuộc ngữ cảnh. Ba nhóm này minh họa các yêu cầu về mô hình hóa thời gian, chi tiết không gian và ngữ cảnh đã nêu cho bài toán phân xưởng.

\section{Mục tiêu nghiên cứu}
\label{sec:objectives}

\begin{enumerate}[label=M\arabic*.,leftmargin=1.2cm]
  \item Thiết kế một temporal neck hậu-backbone kết hợp phụ thuộc dài hạn và biến đổi ngắn hạn với chi phí tuyến tính theo độ dài đoạn video.
  \item Giữ khả năng nạp trọng số tiền huấn luyện của Video Swin bằng cách không sửa cấu trúc bên trong backbone.
  \item Áp dụng temporal neck độc lập trên nhiều tầng đặc trưng để tín hiệu nhỏ không bị mất trước 3D RoIAlign.
  \item Bổ sung ngữ cảnh actor--cảnh bằng ACRN.
  \item Đánh giá trên dữ liệu phân xưởng và AVA~v2.2 với nhiều hạt giống và cách chia theo camera; tách rõ kết quả của đề tài khỏi số liệu tham chiếu.
  \item Duy trì tốc độ xử lý đủ cho giám sát theo thời gian thực.
\end{enumerate}

\section{Câu hỏi nghiên cứu}
\label{sec:rq}

\begin{table}[htbp]
\centering
\caption{Câu hỏi nghiên cứu và bằng chứng cần thu thập}
\label{tab:rq}
\begin{tabularx}{\textwidth}{lXX}
\toprule
\textbf{Mã} & \textbf{Câu hỏi} & \textbf{Bằng chứng}\\
\midrule
RQ1 & Có cần temporal neck tại vị trí hậu-backbone, tiền-RoI không, và một neck tiêu chuẩn (temporal MHSA toàn cục) đã đủ chưa? & Ablation A0 $\rightarrow$ A1 $\rightarrow$ A6 trên hai cách chia dữ liệu\\
RQ2 & Mỗi nhánh (Bi-SSM, W-MHSA), hướng quét và cổng sigmoid đóng góp bao nhiêu? & A2--A6: SSM một chiều $\leftrightarrow$ hai chiều; từng nhánh riêng lẻ; cộng trực tiếp $\leftrightarrow$ cổng\\
RQ3 & Áp dụng khối trên hai tầng đặc trưng (đa tỉ lệ) có tốt hơn một tầng? & A6 $\rightarrow$ A7\\
RQ4 & ACRN có cải thiện khả năng nhận dạng các hành vi phụ thuộc ngữ cảnh actor--cảnh? & A7 $\rightarrow$ A8\\
RQ5 & Kết quả có giữ được trên camera chưa thấy khi huấn luyện và trên bộ dữ liệu công khai? & Cách chia theo camera (S2); ba hạt giống; AVA~v2.2\\
\bottomrule
\end{tabularx}
\end{table}

Mã A0--A8 là các cấu hình trong ma trận ablation (\cref{tab:ablation}). RQ1--RQ4 được trả lời bằng thực nghiệm trên cả hai cách chia dữ liệu; RQ5 được trả lời một phần (\cref{sec:rq-answers}).

\section{Đóng góp và mức độ mới}
\label{sec:novelty}

Đóng góp của luận văn nằm ở mức \emph{thiết kế và tích hợp}. Luận văn không tuyên bố mô hình không gian trạng thái, attention cục bộ, đặc trưng đa tỉ lệ hay ACRN là các khái niệm mới: SSM đã được dùng cho video trong VideoMamba và Video Mamba Suite \cite{li2024videomamba,chen2024videomambasuite}, ACRN được đề xuất từ năm 2018 \cite{sun2018acrn}, còn FPN và MViTv2 đã chứng minh lợi ích của biểu diễn đa tỉ lệ \cite{lin2017fpn,li2022mvitv2}. Điểm mới nằm ở vị trí đặt, cách tham số hóa theo tầng và chuỗi tích hợp cho STAD theo từng người:
\begin{enumerate}[label=(\roman*)]
  \item khối lai hai nhánh: bidirectional selective SSM truyền trạng thái trên toàn bộ trục thời gian với chi phí tuyến tính theo $T$, windowed temporal attention với $W = 5$ giữ trực tiếp các chuyển tiếp cục bộ;
  \item cổng sigmoid học tỷ trọng giữa hai nhánh theo nội dung, theo từng kênh \emph{và} từng bước thời gian, thay vì cộng với hệ số cố định;
  \item vị trí hậu-backbone, tiền-RoI: không phải sửa Video Swin tiền huấn luyện mà vẫn giữ ngữ cảnh cảnh vật trước khi cắt vùng người;
  \item áp dụng độc lập trên hai tầng đặc trưng với tham số riêng, rồi nối multi-scale RoI và ACRN thành pipeline Video Swin $\rightarrow$ Hybrid FGTN-Mamba đa tầng $\rightarrow$ multi-scale RoI $\rightarrow$ ACRN $\rightarrow$ head đa nhãn, cài đặt tự chứa trên MMAction2.
\end{enumerate}

% Tạm bỏ mục 1.8 khỏi bản in; giữ nội dung để có thể khôi phục.
% \section{Phân tầng bằng chứng}
% \label{sec:evidence}
%
% Để người đọc phân biệt rõ nguồn gốc của mỗi con số, luận văn phân tầng bằng chứng theo \cref{tab:evidence}.
%
% \begin{table}[htbp]
% \centering
% \caption{Phân tầng bằng chứng được dùng trong luận văn}
% \label{tab:evidence}
% \begin{tabularx}{\textwidth}{lX}
% \toprule
% \textbf{Mức} & \textbf{Nội dung}\\
% \midrule
% E1 --- Có mã nguồn & Cài đặt phương pháp, cấu hình mô hình và cấu trúc các thành phần trong kho mã\\
% E2 --- Hồ sơ thực nghiệm & Kết quả trên bộ dữ liệu phân xưởng (so sánh sáu mô hình; ablation A0--A8 trên hai cách chia, trung bình ba hạt giống; Precision, Recall, tốc độ, bộ nhớ GPU) và 39,1 mAP trên AVA~v2.2; nhật ký chi tiết không nằm trong kho mã\\
% E3 --- Có mã, chưa có số liệu & Ablation hàm mất mát (A9); AP theo từng hành vi; độ lệch chuẩn giữa các hạt giống; độ trễ đầu--cuối\\
% E4 --- Tham chiếu & Số liệu của A*, STAR, OV-STAD, ST-CLIP và các baseline công bố, trích từ bài báo gốc\\
% \bottomrule
% \end{tabularx}
% \end{table}

\section*{Tiểu kết chương}
\addcontentsline{toc}{section}{Tiểu kết chương}

Chương này đã phát biểu bài toán phát hiện hành động không--thời gian theo từng người cho bốn hành vi mục tiêu trong phân xưởng, chỉ ra hai thang thời gian đối lập của các hành vi và yêu cầu kéo theo đối với toán tử thời gian, đồng thời xác lập độ đo đánh giá, sáu mục tiêu và năm câu hỏi nghiên cứu. \Cref{ch:related} tiếp theo trình bày cơ sở lý thuyết và các họ phương pháp liên quan để làm rõ khoảng trống mà thiết kế đề xuất nhắm tới.
